\textbf{Why tree ensembles dominate.} CICFlowMeter features have heterogeneous scales and heavy tails, and many attacks are separated by sharp thresholds on a few features (destination port, initial window size, header lengths and packet-length maxima; Table~\ref{tab:features}). Axis-aligned tree splits capture such rules directly. The approximate RBF SVM must express the same boundaries through a smooth kernel on standardised features, using 2048 random features and a 500k-row sample. Validation-based tuning cut its false-positive rate from \SvmValFprBefore\% to \SvmValFprAfter\%, but it remains well behind the trees.

\textbf{Why the split matters so much.} \DupPct\% of all flows are exact duplicates. Under the family-stratified split, \EoneOverlap\% of test flows have an identical feature vector in the training set, compared with \EtwoOverlap\% under the capture-file split. Closed-set results are therefore partly a measure of recognition, not generalisation. The large drop in Experiment~2 matches the ``closed world'' concern of Sommer and Paxson~\cite{sommer2010outside}: models learn the attacks and the benign traffic of the days they saw.

\textbf{Value of the anomaly stage.} The autoencoder alone is a weak detector (ROC-AUC \EoneAeAuc, \EoneAeTpr\% of attacks flagged at \EoneAeFpr\% benign FPR), but it is complementary. It flags exactly the rare families the gate misses, and these become \texttt{anomaly\_review} alerts instead of silent misses. The anomaly stage contributes coverage of unusual traffic, not overall accuracy.

\textbf{Feature importance and shortcuts.} The destination port is the single most important Random Forest feature. In CICIDS2017 several attacks target fixed ports (for example FTP and SSH brute forcing, and web attacks on port 80), so this feature can act as a dataset-specific shortcut that would not transfer to networks with different services.
